% =============================================================================
% APPENDIX EV: LIT-VANDENSTEEN
% Eric Van den Steen Research Integration
% =============================================================================
% Category: LIT (Literature Integration)
% Language: English
% Version: 1.0 (January 2026)
%
% Dependencies: CORE-WHO (AAA), CORE-WHEN (V), CORE-AWARE (AU), DOMAIN-ORG
% Related: LIT-TIROLE (JT), LIT-BENABOU (RB), LIT-AKERLOF (GA)
%
% Chapter Linkage:
%   Primary: Chapter 9 (Context), Chapter 10 (Utility Architecture)
%   Secondary: Chapter 5 (Complementarity), Chapter 14 (Behavioral Segments)
% =============================================================================

\chapter{EV LIT-VANDENSTEEN: Organizational Economics \& Corporate Culture}
\label{app:lit-vandensteen}

% -----------------------------------------------------------------------------
% HEADER BLOCK
% -----------------------------------------------------------------------------
\begin{tcolorbox}[colback=blue!5!white, colframe=blue!75!black, title=Appendix Metadata]
\textbf{Code:} EV \\
\textbf{Category:} LIT (Literature Integration) \\
\textbf{Full Name:} LIT-VANDENSTEEN: Eric Van den Steen Research Integration \\
\textbf{Version:} 1.0 (January 2026) \\
\textbf{Papers Covered:} 30 \\
\textbf{Primary Chapters:} 9, 10 \\
\textbf{Key Connections:} CORE-WHEN (V), CORE-AWARE (AU), DOMAIN-ORG, LIT-TIROLE
\end{tcolorbox}

% -----------------------------------------------------------------------------
% CROSS-REFERENCE MAP
% -----------------------------------------------------------------------------
\begin{tcolorbox}[colback=gray!5!white, colframe=gray!75!black, title=Cross-Reference Map]
\textbf{Outgoing References:}
\begin{itemize}[nosep]
    \item CORE-WHO (AAA): Organizational levels (L=3, L=4)
    \item CORE-WHEN (V): Corporate culture as context $\Psi$
    \item CORE-AWARE (AU): Shared beliefs and awareness
    \item CORE-HOW (B): Complementarity in strategy
    \item DOMAIN-ORG: Organizational applications
    \item LIT-TIROLE (JT): Industrial organization foundations
    \item LIT-BENABOU (RB): Motivation and identity
    \item LIT-AKERLOF (GA): Identity economics
\end{itemize}

\textbf{Incoming References:}
\begin{itemize}[nosep]
    \item Chapter 9: Context framework ($\Psi$)
    \item Chapter 10: Utility architecture (IDN dimension)
    \item Chapter 5: Complementarity in organizations
\end{itemize}
\end{tcolorbox}

% =============================================================================
% -----------------------------------------------------------------------------
% APPENDIX SCOPE BOX
% -----------------------------------------------------------------------------

\begin{tcolorbox}[
    colback=orange!5!white,
    colframe=orange!75!black,
    title={\textbf{Appendix Scope: Ziel / In-Scope / Out-of-Scope / Constraints}},
    fonttitle=\bfseries
]

\textbf{Ziel (Objective):}
\begin{itemize}[nosep]
    \item How do shared beliefs emerge, persist, and affect organizational performance?
\end{itemize}

\textbf{In-Scope:}
\begin{itemize}[nosep]
    \item Abstract
    \item Fundamental Question
    \item The Six Van den Steen Axioms
    \item Integration with EBF Framework
    \item Worked Example: Corporate Culture Change
\end{itemize}

\textbf{Out-of-Scope (delegiert an andere Appendices):}
\begin{itemize}[nosep]
    \item REF-GLOSSARY $\rightarrow$ Appendix G
\end{itemize}

\textbf{Constraints (Anwendungsgrenzen):}
\begin{itemize}[nosep]
    \item [TODO: Define when this appendix does NOT apply]
    \item [TODO: Define required prerequisites]
    \item [TODO: Define known limitations]
\end{itemize}

\textbf{Lieferobjekte (Deliverables):}
\begin{itemize}[nosep]
    \item 2 Table(s)
    \item 7 Equation(s)
    \item 12 Axiom(s)
    \item Worked Example(s)
\end{itemize}

\end{tcolorbox}


\section{Abstract}
% =============================================================================

Eric Van den Steen's research program provides foundational theory for understanding
\textbf{corporate culture}, \textbf{shared beliefs}, \textbf{vision}, and
\textbf{organizational strategy} within economic frameworks. His work directly
informs the EBF's treatment of context ($\Psi$), organizational-level utility
(L=3, L=4), and the identity dimension (IDN) of welfare.

This appendix integrates 30 papers spanning 2002--2024, organized around six
core axioms that connect Van den Steen's theoretical contributions to EBF's
complementarity-context framework.

% =============================================================================
\section{Quick Reference}
% =============================================================================

\begin{tcolorbox}[colback=yellow!5!white, colframe=yellow!75!black, title=Key Concepts]
\begin{tabular}{@{}ll@{}}
\textbf{Corporate Culture} & Shared beliefs/values that persist in organizations \\
\textbf{Vision} & Leader's beliefs about optimal strategy \\
\textbf{Homogeneity} & Degree of belief alignment in organizations \\
\textbf{Authority} & Right to make decisions binding on others \\
\textbf{Strategy} & Coherent set of choices creating competitive advantage \\
\textbf{Rational Overoptimism} & Bayesian-rational agents appearing overconfident \\
\end{tabular}
\end{tcolorbox}

% =============================================================================
\section{Fundamental Question}
% =============================================================================

\begin{tcolorbox}[colback=red!5!white, colframe=red!75!black,
    title=The Van den Steen Question]
\centering\Large
\textbf{How do shared beliefs emerge, persist, and affect organizational performance?}
\end{tcolorbox}

Van den Steen's answer: Shared beliefs arise through \textit{sorting} (people
join organizations whose beliefs match theirs) and \textit{learning} (beliefs
converge through shared experience). These beliefs constitute corporate culture,
which affects coordination, motivation, and strategic coherence.

% =============================================================================
\section{The Six Van den Steen Axioms}
% =============================================================================

Van den Steen's contributions can be organized into six core axioms that
integrate with EBF:

\subsection{Axiom UNMAPPED_EV-1: Culture as Shared Beliefs}

\begin{tcolorbox}[colback=blue!5!white, colframe=blue!75!black]
\textbf{Axiom UNMAPPED_EV-1 (Culture Definition):}
\textit{Corporate culture consists of shared beliefs, values, and norms that
emerge endogenously through sorting and learning, and persist across
organizational membership changes.}
\end{tcolorbox}

\textbf{Key Papers:}
\begin{itemize}[nosep]
    \item Van den Steen (2010a): ``Culture Clash: The Costs and Benefits of Homogeneity''
    \item Van den Steen (2010b): ``On the Origin of Shared Beliefs (and Corporate Culture)''
    \item Gibbons \& Van den Steen (2015): ``What Is a Corporate Culture?''
\end{itemize}

\textbf{EBF Integration:}
Corporate culture enters EBF as a component of organizational context $\Psi^{org}$:
\begin{equation}
\Psi^{org} = \{\Psi_{culture}, \Psi_{structure}, \Psi_{incentives}, \Psi_{environment}\}
\end{equation}

Culture affects level salience $\alpha^L(\Psi)$ by making organizational identity
(L=3, L=4) more salient when culture is strong.

\subsection{Axiom UNMAPPED_EV-2: Vision and Leadership}

\begin{tcolorbox}[colback=blue!5!white, colframe=blue!75!black]
\textbf{Axiom UNMAPPED_EV-2 (Vision Effect):}
\textit{A leader's vision---their beliefs about optimal strategy---affects
organizational performance through direction-setting, motivation, and
coordination, even when the vision is ``wrong.''}
\end{tcolorbox}

\textbf{Key Papers:}
\begin{itemize}[nosep]
    \item Van den Steen (2005): ``Organizational Beliefs and Managerial Vision''
    \item Van den Steen (2019): ``The Role of Vision in Leadership and Strategy''
\end{itemize}

\textbf{EBF Integration:}
Vision functions as a focal point (cf. LIT-SCHELLING) that coordinates expectations:
\begin{equation}
WAX^{org} = f(Vision_{clarity}, Belief_{alignment}, \gamma_{coordination})
\end{equation}

Strong vision increases organizational willingness to act by reducing coordination
uncertainty.

\subsection{Axiom UNMAPPED_EV-3: Sorting and Homogeneity}

\begin{tcolorbox}[colback=blue!5!white, colframe=blue!75!black]
\textbf{Axiom UNMAPPED_EV-3 (Sorting Mechanism):}
\textit{Organizations become homogeneous in beliefs through self-selection:
people join organizations whose beliefs match theirs, creating endogenous
corporate culture without explicit socialization.}
\end{tcolorbox}

\textbf{Key Papers:}
\begin{itemize}[nosep]
    \item Van den Steen (2010a): ``Culture Clash''
    \item Van den Steen (2006): ``The Origins of Shared Beliefs''
\end{itemize}

\textbf{EBF Integration:}
Sorting creates identity complementarity within organizations:
\begin{equation}
\gamma^{IDN}_{ij} > 0 \quad \text{when beliefs align through sorting}
\end{equation}

This explains why organizational-level utility (L=3) can exceed the sum of
individual utilities---shared beliefs create emergent value.

\subsection{Axiom UNMAPPED_EV-4: Strategy as Complementarity}

\begin{tcolorbox}[colback=blue!5!white, colframe=blue!75!black]
\textbf{Axiom UNMAPPED_EV-4 (Strategy Coherence):}
\textit{Strategy is a coherent set of interdependent choices, where coherence
derives from complementarity: choices reinforce each other, creating competitive
advantage that is difficult to imitate.}
\end{tcolorbox}

\textbf{Key Papers:}
\begin{itemize}[nosep]
    \item Van den Steen (2017): ``A Formal Theory of Strategy''
    \item Van den Steen (2018): ``Strategy and the Strategist''
    \item Van den Steen (2013): ``The Structure of Strategy''
\end{itemize}

\textbf{EBF Integration:}
This directly connects to EBF's CORE-HOW (B) complementarity framework:
\begin{equation}
U^{strategy} = \sum_i u_i(x_i) + \sum_{i<j} \gamma_{ij} \cdot x_i \cdot x_j
\end{equation}

Strategic complementarity ($\gamma_{ij} > 0$) means the whole exceeds the sum
of parts---a key EBF insight applied to organizational strategy.

\subsection{Axiom UNMAPPED_EV-5: Rational Overoptimism}

\begin{tcolorbox}[colback=blue!5!white, colframe=blue!75!black]
\textbf{Axiom UNMAPPED_EV-5 (Belief Bias):}
\textit{Bayesian-rational agents can appear systematically overconfident and
overoptimistic as an equilibrium outcome, without any cognitive bias---differences
in priors lead to self-serving attribution of success.}
\end{tcolorbox}

\textbf{Key Papers:}
\begin{itemize}[nosep]
    \item Van den Steen (2004): ``Rational Overoptimism (and Other Biases)''
    \item Van den Steen (2011): ``Overconfidence by Bayesian-Rational Agents''
    \item Van den Steen (2003): ``Skill or Luck?''
\end{itemize}

\textbf{EBF Integration:}
This informs CORE-AWARE (AU): apparent biases may reflect heterogeneous beliefs
rather than cognitive limitations:
\begin{equation}
A^{perceived}_{bias} = A^{true} + \Delta_{belief\_heterogeneity}
\end{equation}

Implications for EBF: Be careful attributing ``irrationality'' to agents---apparent
biases may be rational given different priors.

\subsection{Axiom UNMAPPED_EV-6: Authority and Delegation}

\begin{tcolorbox}[colback=blue!5!white, colframe=blue!75!black]
\textbf{Axiom UNMAPPED_EV-6 (Authority Allocation):}
\textit{Authority---the right to make binding decisions---is optimally allocated
based on disagreement: more authority to those whose decisions are less
contestable, enabling commitment and coordination.}
\end{tcolorbox}

\textbf{Key Papers:}
\begin{itemize}[nosep]
    \item Van den Steen (2010): ``Interpersonal Authority in a Theory of the Firm''
    \item Van den Steen (2009): ``Authority versus Persuasion''
    \item Van den Steen (2012): ``Disagreement and the Allocation of Control''
\end{itemize}

\textbf{EBF Integration:}
Authority allocation affects organizational WAX through the $\theta$ threshold:
\begin{equation}
\theta^{action}_{org} = f(Authority_{structure}, Disagreement, Trust)
\end{equation}

Clear authority reduces action thresholds by eliminating coordination uncertainty.

% =============================================================================

% === AUTO-GENERATED ENTRIES (2026-02-22) ===

%%%%%%%%%%%%%%%%%%%%%%%%%%%%%%%%%%%%%%%%%%%%%%%%%%%%%%%%%%%%%%%%%%%%%%%%%%%%%%%
\subsubsection{2002: Essays on the Economics of Heterogeneous Beliefs}
%%%%%%%%%%%%%%%%%%%%%%%%%%%%%%%%%%%%%%%%%%%%%%%%%%%%%%%%%%%%%%%%%%%%%%%%%%%%%%%

\paragraph{Citation.}
Eric Van den Steen (2002). ``Essays on the Economics of Heterogeneous Beliefs.'' \textit{PhD Thesis}.

\paragraph{10C Integration.}
\begin{itemize}[nosep]
  \item \textbf{Domain}: [To be specified]
  \item \textbf{use\_for}: LIT-VANDENSTEEN
\end{itemize}


%%%%%%%%%%%%%%%%%%%%%%%%%%%%%%%%%%%%%%%%%%%%%%%%%%%%%%%%%%%%%%%%%%%%%%%%%%%%%%%
\subsubsection{2007: Evolution of Biases in Rational Agents}
%%%%%%%%%%%%%%%%%%%%%%%%%%%%%%%%%%%%%%%%%%%%%%%%%%%%%%%%%%%%%%%%%%%%%%%%%%%%%%%

\paragraph{Citation.}
Eric Van den Steen (2007). ``Evolution of Biases in Rational Agents.'' \textit{Working Paper}.

\paragraph{10C Integration.}
\begin{itemize}[nosep]
  \item \textbf{Domain}: [To be specified]
  \item \textbf{use\_for}: LIT-VANDENSTEEN
\end{itemize}


%%%%%%%%%%%%%%%%%%%%%%%%%%%%%%%%%%%%%%%%%%%%%%%%%%%%%%%%%%%%%%%%%%%%%%%%%%%%%%%
\subsubsection{2008: On the Origins of Disagreement}
%%%%%%%%%%%%%%%%%%%%%%%%%%%%%%%%%%%%%%%%%%%%%%%%%%%%%%%%%%%%%%%%%%%%%%%%%%%%%%%

\paragraph{Citation.}
Eric Van den Steen (2008). ``On the Origins of Disagreement.'' \textit{Working Paper}.

\paragraph{10C Integration.}
\begin{itemize}[nosep]
  \item \textbf{Domain}: [To be specified]
  \item \textbf{use\_for}: LIT-VANDENSTEEN
\end{itemize}


%%%%%%%%%%%%%%%%%%%%%%%%%%%%%%%%%%%%%%%%%%%%%%%%%%%%%%%%%%%%%%%%%%%%%%%%%%%%%%%
\subsubsection{2011: Means versus Ends in Opaque Institutional Environments}
%%%%%%%%%%%%%%%%%%%%%%%%%%%%%%%%%%%%%%%%%%%%%%%%%%%%%%%%%%%%%%%%%%%%%%%%%%%%%%%

\paragraph{Citation.}
Eric Van den Steen (2011). ``Means versus Ends in Opaque Institutional Environments.'' \textit{Working Paper}.

\paragraph{10C Integration.}
\begin{itemize}[nosep]
  \item \textbf{Domain}: [To be specified]
  \item \textbf{use\_for}: LIT-VANDENSTEEN
\end{itemize}


%%%%%%%%%%%%%%%%%%%%%%%%%%%%%%%%%%%%%%%%%%%%%%%%%%%%%%%%%%%%%%%%%%%%%%%%%%%%%%%
\subsubsection{2014: Why Do Firms Have 'Purpose'? The Firm's Role as Carrier of I...}
%%%%%%%%%%%%%%%%%%%%%%%%%%%%%%%%%%%%%%%%%%%%%%%%%%%%%%%%%%%%%%%%%%%%%%%%%%%%%%%

\paragraph{Citation.}
Eric Van den Steen (2014). ``Why Do Firms Have 'Purpose'? The Firm's Role as Carrier of Identity and Reputation.'' \textit{Working Paper}.

\paragraph{10C Integration.}
\begin{itemize}[nosep]
  \item \textbf{Domain}: [To be specified]
  \item \textbf{use\_for}: LIT-VANDENSTEEN
\end{itemize}


%%%%%%%%%%%%%%%%%%%%%%%%%%%%%%%%%%%%%%%%%%%%%%%%%%%%%%%%%%%%%%%%%%%%%%%%%%%%%%%
\subsubsection{2014: Tesla Motors: Strategy and Business Model}
%%%%%%%%%%%%%%%%%%%%%%%%%%%%%%%%%%%%%%%%%%%%%%%%%%%%%%%%%%%%%%%%%%%%%%%%%%%%%%%

\paragraph{Citation.}
Eric Van den Steen (2014). ``Tesla Motors: Strategy and Business Model.'' \textit{Harvard Business School Case}.

\paragraph{10C Integration.}
\begin{itemize}[nosep]
  \item \textbf{Domain}: [To be specified]
  \item \textbf{use\_for}: LIT-VANDENSTEEN
\end{itemize}


%%%%%%%%%%%%%%%%%%%%%%%%%%%%%%%%%%%%%%%%%%%%%%%%%%%%%%%%%%%%%%%%%%%%%%%%%%%%%%%
\subsubsection{2015: Coordination in Organizations: A Game-Theoretic Perspective}
%%%%%%%%%%%%%%%%%%%%%%%%%%%%%%%%%%%%%%%%%%%%%%%%%%%%%%%%%%%%%%%%%%%%%%%%%%%%%%%

\paragraph{Citation.}
Eric Van den Steen (2015). ``Coordination in Organizations: A Game-Theoretic Perspective.'' \textit{Working Paper}.

\paragraph{10C Integration.}
\begin{itemize}[nosep]
  \item \textbf{Domain}: [To be specified]
  \item \textbf{use\_for}: LIT-VANDENSTEEN
\end{itemize}


%%%%%%%%%%%%%%%%%%%%%%%%%%%%%%%%%%%%%%%%%%%%%%%%%%%%%%%%%%%%%%%%%%%%%%%%%%%%%%%
\subsubsection{2016: OpenAI: Strategy and Vision}
%%%%%%%%%%%%%%%%%%%%%%%%%%%%%%%%%%%%%%%%%%%%%%%%%%%%%%%%%%%%%%%%%%%%%%%%%%%%%%%

\paragraph{Citation.}
Eric Van den Steen (2016). ``OpenAI: Strategy and Vision.'' \textit{Harvard Business School Case}.

\paragraph{10C Integration.}
\begin{itemize}[nosep]
  \item \textbf{Domain}: [To be specified]
  \item \textbf{use\_for}: LIT-VANDENSTEEN
\end{itemize}


%%%%%%%%%%%%%%%%%%%%%%%%%%%%%%%%%%%%%%%%%%%%%%%%%%%%%%%%%%%%%%%%%%%%%%%%%%%%%%%
\subsubsection{2016: Strategy Mistakes: How to Avoid the Most Common Errors}
%%%%%%%%%%%%%%%%%%%%%%%%%%%%%%%%%%%%%%%%%%%%%%%%%%%%%%%%%%%%%%%%%%%%%%%%%%%%%%%

\paragraph{Citation.}
Eric Van den Steen (2016). ``Strategy Mistakes: How to Avoid the Most Common Errors.'' \textit{Harvard Business Review}.

\paragraph{10C Integration.}
\begin{itemize}[nosep]
  \item \textbf{Domain}: [To be specified]
  \item \textbf{use\_for}: LIT-VANDENSTEEN
\end{itemize}


%%%%%%%%%%%%%%%%%%%%%%%%%%%%%%%%%%%%%%%%%%%%%%%%%%%%%%%%%%%%%%%%%%%%%%%%%%%%%%%
\subsubsection{2020: Does Corporate Culture Matter for Performance?}
%%%%%%%%%%%%%%%%%%%%%%%%%%%%%%%%%%%%%%%%%%%%%%%%%%%%%%%%%%%%%%%%%%%%%%%%%%%%%%%

\paragraph{Citation.}
Eric Van den Steen (2020). ``Does Corporate Culture Matter for Performance?.'' \textit{Working Paper}.

\paragraph{10C Integration.}
\begin{itemize}[nosep]
  \item \textbf{Domain}: [To be specified]
  \item \textbf{use\_for}: LIT-VANDENSTEEN
\end{itemize}


%%%%%%%%%%%%%%%%%%%%%%%%%%%%%%%%%%%%%%%%%%%%%%%%%%%%%%%%%%%%%%%%%%%%%%%%%%%%%%%
\subsubsection{2021: Making Strategy: The Role of Strategy in Practice}
%%%%%%%%%%%%%%%%%%%%%%%%%%%%%%%%%%%%%%%%%%%%%%%%%%%%%%%%%%%%%%%%%%%%%%%%%%%%%%%

\paragraph{Citation.}
Eric Van den Steen (2021). ``Making Strategy: The Role of Strategy in Practice.'' \textit{Harvard Business School Case}.

\paragraph{10C Integration.}
\begin{itemize}[nosep]
  \item \textbf{Domain}: [To be specified]
  \item \textbf{use\_for}: LIT-VANDENSTEEN
\end{itemize}


%%%%%%%%%%%%%%%%%%%%%%%%%%%%%%%%%%%%%%%%%%%%%%%%%%%%%%%%%%%%%%%%%%%%%%%%%%%%%%%
\subsubsection{2022: Strategy Under Uncertainty}
%%%%%%%%%%%%%%%%%%%%%%%%%%%%%%%%%%%%%%%%%%%%%%%%%%%%%%%%%%%%%%%%%%%%%%%%%%%%%%%

\paragraph{Citation.}
Eric Van den Steen (2022). ``Strategy Under Uncertainty.'' \textit{Working Paper}.

\paragraph{10C Integration.}
\begin{itemize}[nosep]
  \item \textbf{Domain}: [To be specified]
  \item \textbf{use\_for}: LIT-VANDENSTEEN
\end{itemize}


%%%%%%%%%%%%%%%%%%%%%%%%%%%%%%%%%%%%%%%%%%%%%%%%%%%%%%%%%%%%%%%%%%%%%%%%%%%%%%%
\subsubsection{2023: The Nature of Strategic Choice}
%%%%%%%%%%%%%%%%%%%%%%%%%%%%%%%%%%%%%%%%%%%%%%%%%%%%%%%%%%%%%%%%%%%%%%%%%%%%%%%

\paragraph{Citation.}
Eric Van den Steen (2023). ``The Nature of Strategic Choice.'' \textit{Working Paper}.

\paragraph{10C Integration.}
\begin{itemize}[nosep]
  \item \textbf{Domain}: [To be specified]
  \item \textbf{use\_for}: LIT-VANDENSTEEN
\end{itemize}


%%%%%%%%%%%%%%%%%%%%%%%%%%%%%%%%%%%%%%%%%%%%%%%%%%%%%%%%%%%%%%%%%%%%%%%%%%%%%%%
\subsubsection{2024: The Dynamics of Corporate Culture}
%%%%%%%%%%%%%%%%%%%%%%%%%%%%%%%%%%%%%%%%%%%%%%%%%%%%%%%%%%%%%%%%%%%%%%%%%%%%%%%

\paragraph{Citation.}
Eric Van den Steen (2024). ``The Dynamics of Corporate Culture.'' \textit{Working Paper}.

\paragraph{10C Integration.}
\begin{itemize}[nosep]
  \item \textbf{Domain}: [To be specified]
  \item \textbf{use\_for}: LIT-VANDENSTEEN
\end{itemize}

% === END AUTO-GENERATED ENTRIES ===
\section{Summary of Key Papers}
% =============================================================================

\begin{table}[htbp]
\centering
\caption{Van den Steen Papers: Evidence Tiers and EBF Integration}
\label{tab:vandensteen-papers}
\renewcommand{\arraystretch}{1.2}
\small
\begin{tabular}{@{}p{4cm}cp{5cm}c@{}}
\toprule
\textbf{Paper} & \textbf{Year} & \textbf{Key Contribution} & \textbf{Tier} \\
\midrule
Culture Clash & 2010 & Costs/benefits of homogeneity & 5 \\
Interpersonal Authority & 2010 & Theory of firm boundaries & 5 \\
Organizational Beliefs & 2005 & Vision and leadership & 5 \\
Formal Theory of Strategy & 2017 & Strategy as complementarity & 5 \\
Rational Overoptimism & 2004 & Biases from heterogeneous beliefs & 5 \\
Origin of Shared Beliefs & 2010 & Sorting creates culture & 5 \\
Authority vs Persuasion & 2009 & When to use each & 5 \\
Strategy and Strategist & 2018 & Who formulates matters & 5 \\
\bottomrule
\end{tabular}
\end{table}

% =============================================================================

%%%%%%%%%%%%%%%%%%%%%%%%%%%%%%%%%%%%%%%%%%%%%%%%%%%%%%%%%%%%%%%%%%%%%%%%%%%%%%%
\section{Key Results}
\label{sec:results}
%%%%%%%%%%%%%%%%%%%%%%%%%%%%%%%%%%%%%%%%%%%%%%%%%%%%%%%%%%%%%%%%%%%%%%%%%%%%%%%

The analysis yields the following key results:

\begin{enumerate}
    \item \textbf{Result 1:} [TODO: Describe main finding]
    \item \textbf{Result 2:} [TODO: Describe secondary finding]
    \item \textbf{Result 3:} [TODO: Describe implication]
\end{enumerate}

\section{Integration with EBF Framework}
% =============================================================================

\subsection{Connection to 10C CORE Questions}

\begin{table}[htbp]
\centering
\caption{Van den Steen $\rightarrow$ 10C Mapping}
\label{tab:vandensteen-9c}
\renewcommand{\arraystretch}{1.3}
\begin{tabular}{@{}lll@{}}
\toprule
\textbf{10C Question} & \textbf{Van den Steen Contribution} & \textbf{Mechanism} \\
\midrule
WHO (AAA) & Organizational levels L=3, L=4 & Firms as welfare-bearing entities \\
WHAT (C) & IDN dimension & Corporate identity, purpose \\
HOW (B) & Strategic complementarity & $\gamma > 0$ in strategy \\
WHEN (V) & Culture as context & $\Psi_{culture}$ modulates behavior \\
AWARE (AU) & Shared beliefs & Collective awareness function \\
READY (AV) & Authority reduces $\theta$ & Clear authority enables action \\
\bottomrule
\end{tabular}
\end{table}

\subsection{Theoretical Synthesis}

Van den Steen's work provides microfoundations for organizational-level phenomena
in EBF:

\begin{enumerate}
    \item \textbf{Culture = Shared $\Psi$}: Corporate culture is the shared
    context that organizational members experience, affecting their level
    salience and decision-making.

    \item \textbf{Vision = Focal Awareness}: Leadership vision coordinates
    what people are aware of ($A$) and willing to act on ($WAX$).

    \item \textbf{Strategy = Complementarity System}: Strategic choices exhibit
    $\gamma > 0$, creating emergent value beyond individual choices.

    \item \textbf{Sorting = Endogenous Homogeneity}: Organizations become
    belief-homogeneous through selection, not socialization.
\end{enumerate}

% =============================================================================
\section{Worked Example: Corporate Culture Change}
% =============================================================================

\begin{tcolorbox}[colback=green!5!white, colframe=green!75!black,
    title=Example: Culture Transformation at ``TechCorp'']

\textbf{Situation:} TechCorp wants to shift from hierarchical to agile culture.

\textbf{Van den Steen Analysis:}
\begin{enumerate}
    \item \textbf{UNMAPPED_EV-1 (Culture as Beliefs):} Current culture = shared beliefs
    about ``how we work here.'' Must change beliefs, not just practices.

    \item \textbf{UNMAPPED_EV-3 (Sorting):} New hires will self-select based on perceived
    culture. To change culture, must signal new beliefs to attract different people.

    \item \textbf{UNMAPPED_EV-2 (Vision):} Leaders must articulate clear vision of new
    culture to coordinate expectations.

    \item \textbf{UNMAPPED_EV-6 (Authority):} Agile requires different authority allocation---
    more delegation, less hierarchy.
\end{enumerate}

\textbf{EBF Prediction:}
\begin{equation}
P(culture\_change) = f(\underbrace{Vision_{clarity}}_{\text{UNMAPPED_EV-2}},
\underbrace{Turnover_{rate}}_{\text{UNMAPPED_EV-3}},
\underbrace{Authority_{reallocation}}_{\text{UNMAPPED_EV-6}})
\end{equation}

Culture change requires all three: clear vision, sufficient turnover to enable
sorting, and aligned authority structures.

\end{tcolorbox}

% =============================================================================
\section{Critical Foundations}
% =============================================================================

\begin{enumerate}
    \item \textbf{Methodological:} Van den Steen uses formal game-theoretic
    models, providing rigorous microfoundations (Evidence Tier 5).

    \item \textbf{Empirical Gap:} Limited direct empirical tests of key
    predictions (culture-performance link, sorting mechanisms).

    \item \textbf{Complementarity Link:} Strategy-as-complementarity directly
    supports EBF's CORE-HOW framework.

    \item \textbf{Rationality Assumption:} Maintains Bayesian rationality while
    generating ``behavioral'' predictions---important bridge for EBF.

    \item \textbf{Organizational Focus:} Primarily L=3 (organization) level,
    with implications for L=4 (network) through inter-firm culture.
\end{enumerate}

% =============================================================================
\section{Open Issues}
% =============================================================================

\begin{itemize}
    \item \textbf{Empirical Validation:} How to measure ``shared beliefs''
    directly? Current proxies (surveys, artifacts) are noisy.

    \item \textbf{Culture Dynamics:} How fast can culture change? Van den Steen's
    sorting mechanism suggests slow change bounded by turnover.

    \item \textbf{Multi-Level Integration:} How do individual beliefs (L=0)
    aggregate to organizational culture (L=3)?

    \item \textbf{Cross-Cultural Application:} Does the theory apply equally
    across national cultures with different baseline beliefs?
\end{itemize}

% =============================================================================
\section{Glossary}
% =============================================================================

For complete definitions, see \textbf{Appendix G (REF-GLOSSARY)}.

\begin{description}
    \item[Corporate Culture] Shared beliefs, values, and norms within an organization
    \item[Vision] Leader's beliefs about optimal organizational strategy
    \item[Sorting] Self-selection mechanism creating belief homogeneity
    \item[Authority] Right to make decisions binding on others
    \item[Strategic Complementarity] Choices that reinforce each other ($\gamma > 0$)
\end{description}

% =============================================================================
\section{References}
% =============================================================================

\nocite{bcm_master}

Van den Steen's 30 papers (2002--2024) are catalogued in the EBF paper database
under category \texttt{LIT-VANDENSTEEN}. Key journals include:
\begin{itemize}[nosep]
    \item American Economic Review (2)
    \item Management Science (4)
    \item Journal of Law, Economics, and Organization (2)
    \item RAND Journal of Economics (1)
    \item Harvard Business School Working Papers (12)
    \item Harvard Business School Cases (3)
\end{itemize}

\end{document}
